% =========================================================
% PARTE 3: REGRESIÓN LINEAL SIMPLE Y PRUEBAS DE SUPUESTOS
% =========================================================
\section{Parte 3: Modelo de Regresión Lineal Simple}

\subsection{Ecuación del Modelo Teórico y Estimadores OLS}

El modelo probabilístico de regresión lineal simple responde a la formulación:
\begin{equation}
	Y_i = a + b \cdot X_i + e_i, \quad e_i \sim \mathcal{N}(0, \sigma^2)
\end{equation}

Los estimadores de Mínimos Cuadrados Ordinarios (OLS) para la intersección ($a$) y la pendiente ($b$) se calculan mediante:

\begin{statbox}{Formulación de Estimadores OLS}
	\begin{equation}
		b = \frac{\sum_{i=1}^{n} (X_i - \bar{X})(Y_i - \bar{Y})}{\sum_{i=1}^{n} (X_i - \bar{X})^2}, \quad a = \bar{Y} - b \bar{X}
	\end{equation}
\end{statbox}

Se evalúa específicamente el \textbf{Modelo 2 (Tiempo para Llenar Combi vs Alumnos Esperando)}:
\begin{equation}
	\text{minutos\_para\_llenar} = a + b \times \text{alumnos\_esperando}
\end{equation}

\subsection{Contraste de Hipótesis Estadísticas}

Para evaluar la significancia del coeficiente de la pendiente ($b$), se plantea el siguiente contraste de hipótesis bilateral:

\begin{table}[H]
	\centering
	\caption{Planteamiento de Hipótesis Estadísticas sobre la Pendiente}
	\label{tab:hipotesis_modelo}
	\begin{tabularx}{\textwidth}{>{\bfseries\color{purpleUPP}}l p{4cm} X}
		\toprule
		Hipótesis & Formulación & Interpretación Contextual \\
		\midrule
		Hipótesis Nula ($H_0$) & $H_0: b = 0$ & El número de alumnos en espera \textbf{no influye} en el tiempo restante para llenar la unidad. \\
		\addlinespace[0.5em]
		Hipótesis Alternativa ($H_1$) & $H_1: b \neq 0$ & El número de alumnos en espera \textbf{reduce significativamente} el tiempo necesario para llenar la unidad. \\
		\bottomrule
	\end{tabularx}
\end{table}

\subsection{Verificación de Supuestos del Modelo OLS}

Para validar la inferencia del modelo, se analizan los residuos $e_i = Y_i - \hat{Y}_i$:

\begin{enumerate}
	\item \textbf{Normalidad de Residuos (Prueba de Shapiro-Wilk):}
	Evalúa si los errores provienen de una distribución normal. Se acepta el supuesto si $p_{\text{Shapiro}} > 0.05$.
	\item \textbf{Homocedasticidad (Prueba de Breusch-Pagan):}
	Evalúa si la varianza de los residuos es constante ($\text{Var}(e_i) = \sigma^2$). Se acepta el supuesto si $p_{\text{BP}} > 0.05$.
	\item \textbf{Independencia de Errores (Estadístico de Durbin-Watson):}
	Evalúa la ausencia de autocorrelación. Se exige que el estadístico $DW$ se sitúe en el intervalo $[1.5, 2.5]$.
\end{enumerate}

\begin{figure}[H]
	\centering
	\IfFileExists{img/fig1_regresion_lineal.png}{%
		\includegraphics[width=0.85\textwidth]{img/fig1_regresion_lineal.png}%
	}{%
		\vspace{2em}
		\begin{tcolorbox}[colback=purpleUPP!5!white,colframe=purpleUPP,title={Gráfico en Espera}]
			\centering \textit{La gráfica ajustada ($Y = a + b X$) se generará e incluirá automáticamente al ejecutar el script \texttt{02\_train\_regression.py}.}
		\end{tcolorbox}
		\vspace{2em}
	}%
	\caption{Gráfico de Dispersión y Recta de Regresión Ajustada ($Y = a + b X$).}
	\label{fig:regresion_fit}
\end{figure}

\newpage
